\section{OctoML: de código abierto a startup}

\subsection{Motivación para emprender}

La fundación de OctoML fue muy natural: varios estudiantes y profesores del proyecto TVM vieron una oportunidad de comercialización y decidieron hacer una startup. La motivación central de Chen Tianqi (陈天奇) era «que la comunidad de código abierto pudiera tener éxito»: hacer una startup por el código abierto, no hacer código abierto por la startup.

Durante el primer año participó de tiempo completo, y después pasó a part-time, encargándose principalmente de la dirección técnica y de la investigación de la próxima generación de tecnología. Él tenía «muy claro que no quería ser CEO», porque «tienes que renunciar a mucha tecnología», y su passion estaba en hacer buena tecnología.

\subsection{El cambio de modelo de negocio}

El modelo de negocio de OctoML pasó por un cambio significativo:

\begin{enumerate}
    \item \textbf{Etapa temprana (2019--2022)}: mediante la compilación de aprendizaje automático, ayudaban a los fabricantes de hardware y a los clientes a desplegar modelos con mayor rapidez en distintos tipos de hardware. Los clientes potenciales incluían fabricantes de chips y empresas que necesitaban desplegar ML.
    \item \textbf{Etapa posterior (2023--2024)}: se transformaron en un proveedor de servicios de API de LLM (similar a Fireworks y Together), ofreciendo directamente endpoints de modelos como Llama, con cobro por llamada a la API. La tecnología de compilación de aprendizaje automático pasó a segundo plano, convirtiéndose en tecnología de soporte que ayudaba a reducir el costo de inference.
    \item \textbf{Finales de 2024}: fueron adquiridos por NVIDIA, y el equipo quedó a cargo de la dirección del AI compiler de NVIDIA.
\end{enumerate}

El presentador insistió en que el mercado de API tenía una competencia intensa, en la que todos competían por tener el Llama inference cost más bajo, y preguntó cómo respondía OctoML ante eso. Chen Tianqi reconoció que la competencia era intensa, pero señaló que aún había espacio para diferenciarse en la forma del producto, por ejemplo, dando soporte a multi-LoRA inference. También admitió con franqueza que la adquisición por NVIDIA no fue una ruta de salida planeada de antemano: «Nunca pensé en esa pregunta; en ese momento solo sentía que bastaba con hacer bien una cosa y ver qué pasaba.»

\begin{knowledgebox}{El precedente de la experiencia emprendedora del mentor}
Antes de unirse a OctoML, Chen Tianqi ya había sido testigo de la evolución de Turi (anteriormente GraphLab/Dato), la empresa fundada por su mentor Carlos Guestrin: partió como software de procesamiento de grafos a gran escala en CMU y, para cuando fue adquirida, ya se había transformado por completo en una plataforma de Data Science. «Así que, al empezar a hacer esto, ya tenía la preparación mental de que lo que empiezas a hacer y what you end up doing sin duda serán distintos.» Databricks partió de Spark, y ahora Spark tampoco es ya su parte más importante. Esta es una regularidad general de las startups de código abierto.
\end{knowledgebox}

\begin{warningbox}{La relación sutil entre el código abierto y el negocio}
Chen Tianqi admitió que hacer una startup a partir de un proyecto de código abierto enfrenta el problema de que «el código abierto es tan bueno que nadie paga». Pero su postura es: (1) el código abierto es tu punto de partida, pero casi nunca es tu punto de llegada; el negocio central de Databricks tampoco es ya Spark; (2) lo que toda startup piensa hacer al inicio y lo que termina haciendo «sin duda serán distintos»; (3) una vez que eliges el código abierto, hay que aceptar tanto sus ventajas (retroalimentación de la comunidad, influencia) como sus retos (el diseño del modelo de negocio).
\end{warningbox}

\subsection{El Lesson más grande del emprendimiento}

Cinco años de emprendimiento le dieron a Chen Tianqi importantes conocimientos más allá de lo técnico:

\begin{itemize}
    \item \textbf{Human factor is very important}: «El lesson más grande de hacer una startup es que muchas veces no solo hay factores técnicos. Cómo comunicarte con la gente, cómo desempeñar un role de leadership dentro de un equipo grande: creo que esto es sumamente importante.»
    \item \textbf{Hay que reinvent yourself}: no puedes aferrarte sin cambios a tu ruta técnica original. «Si hace dos años hubiéramos seguido optimizando ResNet, creo que al final habríamos perdido de forma muy contundente.» En el campo de la IA, que cambia con rapidez, la velocidad y el valor para transformarse técnicamente son fundamentales.
    \item \textbf{La excellence técnica nunca es un error}: «Esta opinión no la comparte todo el mundo. Pero, desde mi punto de vista, muchas veces hacer tecnología excelente nunca es un error. Porque in the end, you have to have something to differentiate yourself.»
    \item \textbf{Idealismo + realismo}: «Tengo un poco de idealismo. Pero esta experiencia me hace sentir que soy más realista: sé cómo, por medios realistas, lograr mejor mis ideales.»
    \item \textbf{Product Market Fit}: hay que pensar a fondo a quién y qué problema resuelve realmente tu producto. Pero, al mismo tiempo, no puedes fijarte solo en el PMF de corto plazo e ignorar el valor técnico de largo plazo.
    \item \textbf{Tener claro tu propio rol}: «Tengo muy claro que no quiero ser CEO. Porque ser CEO es sumamente difícil, tienes que renunciar a mucha tecnología. Mi passion está en hacer buena tecnología.» Conocer en qué eres bueno y qué amas, y luego encontrar al co-founder adecuado que te complemente.
\end{itemize}

\begin{warningbox}{La relación sutil entre el código abierto y el negocio}
Chen Tianqi admitió que hacer una startup a partir de un proyecto de código abierto enfrenta el problema de que «el código abierto es tan bueno que nadie paga» (el presentador puso el ejemplo de Anyscale y Ray). Pero su postura es: (1) el código abierto es tu punto de partida, pero casi nunca es tu punto de llegada; el negocio central de Databricks tampoco es ya Spark; (2) lo que toda startup piensa hacer al inicio y lo que termina haciendo «sin duda serán distintos»; (3) una vez que eliges el código abierto, hay que aceptar tanto sus ventajas (retroalimentación de la comunidad, influencia, confianza de los usuarios) como sus retos (el modelo de negocio necesita diferenciarse). Hu Yuanming (胡渊鸣), invitado del episodio anterior, con su Taichi (太极) también enfrentó un problema similar: «Mucha gente dice que Taichi es muy bueno, pero no está dispuesta a pagar.» Este es un reto general para el emprendimiento basado en código abierto.
\end{warningbox}

\subsection{Resumen de la sección}

Los cinco años de trayectoria de OctoML muestran la ruta típica de comercialización de un proyecto académico de código abierto: de estar impulsado por la tecnología a estar impulsado por el producto, de herramienta de compilación a servicio de API, hasta terminar en una adquisición estratégica. En este proceso, Chen Tianqi aprendió dimensiones importantes más allá de lo técnico: el factor humano, el ajuste entre producto y mercado, y la necesidad de una autorrenovación constante.
